%%%%%%%%%%%%%%%%%%%%%%%%%%%%%%%%%%%%%%%%%%%%%%%%%%%%%%%%%%%%%%%%%%%%%%%%%%%%
% AGUJournalTemplate.tex: this template file is for articles formatted with LaTeX
%
% This file includes commands and instructions
% given in the order necessary to produce a final output that will
% satisfy AGU requirements, including customized APA reference formatting.
%
% You may copy this file and give it your
% article name, and enter your text.
%
% guidelines and troubleshooting are here: 

%% To submit your paper:
\documentclass[draft]{agujournal2019}
\usepackage{url} %this package should fix any errors with URLs in refs.
\usepackage{lineno}
\usepackage[inline]{trackchanges} %for better track changes. finalnew option will compile document with changes incorporated.
\usepackage{soul}
\linenumbers
%%%%%%%
% As of 2018 we recommend use of the TrackChanges package to mark revisions.
% The trackchanges package adds five new LaTeX commands:
%
%  \note[editor]{The note}
%  \annote[editor]{Text to annotate}{The note}
%  \add[editor]{Text to add}
%  \remove[editor]{Text to remove}
%  \change[editor]{Text to remove}{Text to add}
%
% complete documentation is here: http://trackchanges.sourceforge.net/
%%%%%%%

\draftfalse

%% Enter journal name below.
%% Choose from this list of Journals:
%
% JGR: Atmospheres
% JGR: Biogeosciences
% JGR: Earth Surface
% JGR: Oceans
% JGR: Planets
% JGR: Solid Earth
% JGR: Space Physics
% Global Biogeochemical Cycles
% Geophysical Research Letters
% Paleoceanography and Paleoclimatology
% Radio Science
% Reviews of Geophysics
% Tectonics
% Space Weather
% Water Resources Research
% Geochemistry, Geophysics, Geosystems
% Journal of Advances in Modeling Earth Systems (JAMES)
% Earth's Future
% Earth and Space Science
% Geohealth
%
% ie, \journalname{Water Resources Research}

\journalname{Earth's Future}


\begin{document}

%%%%%%%%%%%%%%%%%%%%%%%%%%%%%%%%%%%%%%%%%%%%%%%
%  TITLE
%
% (A title should be specific, informative, and brief. Use
% abbreviations only if they are defined in the abstract. Titles that
% start with general keywords then specific terms are optimized in
% searches)
%
%%%%%%%%%%%%%%%%%%%%%%%%%%%%%%%%%%%%%%%%%%%%%%%

% Example: \title{This is a test title}

\title{=enter title here=}

%%%%%%%%%%%%%%%%%%%%%%%%%%%%%%%%%%%%%%%%%%%%%%%
%
%  AUTHORS AND AFFILIATIONS
%
%%%%%%%%%%%%%%%%%%%%%%%%%%%%%%%%%%%%%%%%%%%%%%%

% Authors are individuals who have significantly contributed to the
% research and preparation of the article. Group authors are allowed, if
% each author in the group is separately identified in an appendix.)

% List authors by first name or initial followed by last name and
% separated by commas. Use \affil{} to number affiliations, and
% \thanks{} for author notes.
% Additional author notes should be indicated with \thanks{} (for
% example, for current addresses).

% Example: \authors{A. B. Author\affil{1}\thanks{Current address, Antartica}, B. C. Author\affil{2,3}, and D. E.
% Author\affil{3,4}\thanks{Also funded by Monsanto.}}

\authors{=list all authors here=}


% \affiliation{1}{First Affiliation}
% \affiliation{2}{Second Affiliation}
% \affiliation{3}{Third Affiliation}
% \affiliation{4}{Fourth Affiliation}

\affiliation{=number=}{=Affiliation Address=}
%(repeat as many times as is necessary)


% Corresponding author mailing address and e-mail address:

% (include name and email addresses of the corresponding author.  More
% than one corresponding author is allowed in this LaTeX file and for
% publication; but only one corresponding author is allowed in our
% editorial system.)

% Example: \correspondingauthor{First and Last Name}{email@address.edu}

\correspondingauthor{=name=}{=email address=}



%%%%%%%%%%%%%%%%%%%%%%%%%%%%%%%%%%%%%%%%%%%%%%%
% KEY POINTS
%%%%%%%%%%%%%%%%%%%%%%%%%%%%%%%%%%%%%%%%%%%%%%%
%  List up to three key points (at least one is required)
%  Key Points summarize the main points and conclusions of the article
%  Each must be 140 characters or fewer with no special characters or punctuation and must be complete sentences

% Example:
% \begin{keypoints}
% \item	List up to three key points (at least one is required)
% \item	Key Points summarize the main points and conclusions of the article
% \item	Each must be 140 characters or fewer with no special characters or punctuation and must be complete sentences
% \end{keypoints}

\begin{keypoints}
\item enter point 1 here
\item enter point 2 here
\item enter point 3 here
\end{keypoints}

%%%%%%%%%%%%%%%%%%%%%%%%%%%%%%%%%%%%%%%%%%%%%%%
%
%  ABSTRACT and PLAIN LANGUAGE SUMMARY
%
% A good Abstract will begin with a short description of the problem
% being addressed, briefly describe the new data or analyses, then
% briefly states the main conclusion(s) and how they are supported and
% uncertainties.

% The Plain Language Summary should be written for a broad audience,
% including journalists and the science-interested public, that will not have 
% a background in your field.
%
% A Plain Language Summary is required in GRL, JGR: Planets, JGR: Biogeosciences,
% JGR: Oceans, G-Cubed, Reviews of Geophysics, and JAMES.
% see http://sharingscience.agu.org/creating-plain-language-summary/)
%
%%%%%%%%%%%%%%%%%%%%%%%%%%%%%%%%%%%%%%%%%%%%%%%

%% \begin{abstract} starts the second page

\begin{abstract}
[ enter your Abstract here ]
\end{abstract}

\section*{Plain Language Summary}
Enter your Plain Language Summary here or delete this section.
Here are instructions on writing a Plain Language Summary: 
https://www.agu.org/Share-and-Advocate/Share/Community/Plain-language-summary


%%%%%%%%%%%%%%%%%%%%%%%%%%%%%%%%%%%%%%%%%%%%%%%
%
%  BODY TEXT
%
%%%%%%%%%%%%%%%%%%%%%%%%%%%%%%%%%%%%%%%%%%%%%%%

%%% Suggested section heads:
% \section{Introduction}
%
% The main text should start with an introduction. Except for short
% manuscripts (such as comments and replies), the text should be divided
% into sections, each with its own heading.

% Headings should be sentence fragments and do not begin with a
% lowercase letter or number. Examples of good headings are:

% \section{Materials and Methods}
% Here is text on Materials and Methods.
%
% \subsection{A descriptive heading about methods}
% More about Methods.
%
% \section{Data} (Or section title might be a descriptive heading about data)
%
% \section{Results} (Or section title might be a descriptive heading about the
% results)
%
% \section{Conclusions}


\section{Introduction}

Fecal indicator bacteria (FIB), and \textit{Escherichia coli} in particular, are the primary tool used worldwide to judge whether fresh waters are safe from waterborne pathogens. The public-health stakes are large: inadequate water, sanitation, and hygiene remain responsible for a substantial global disease burden \cite{prssustn2014burden,prssstn2019burden,wolf2023burden}, fecal contamination of drinking water is widespread across low- and middle-income settings \cite{bain2014fecal,bain2014global}, and contact with fecally contaminated recreational water is an established cause of gastrointestinal and other illness \cite{gerdes2023estimating,wade2022health,zmirou2003risks}. Because directly enumerating the many waterborne pathogens is impractical for routine monitoring, agencies instead measure culturable indicator organisms whose presence signals fecal contamination. A long line of epidemiological studies has tied indicator concentrations to swimmer illness and established \textit{E. coli} and intestinal enterococci as the standard indicators for fresh and marine waters respectively \citeA{cabelli1982swimmingassociated}; subsequent reviews and cohort studies refined these relationships \cite{prss1998review,wade2003environmental,wade2005rapidly,wade2010rapidly}, and \textit{E. coli} is widely regarded as the most reliable bacterial indicator of fecal contamination in fresh water \cite{edberg2000escherichia,boehm2014enterococci,korajkic2018relationships}.

Despite this central role, the global record that underpins freshwater safety decisions has never been assembled and characterized at a harmonized, global scale. Monitoring data are dispersed across dozens of national and supranational portals with incompatible schemas, units, and indicator vocabularies, and large-sample syntheses to date have emphasized water chemistry or remote-sensing-amenable variables rather than fecal indicators \cite{read2017water,virro2021grqa,hartmann2014brief}. As a result, two structural weaknesses of the monitoring enterprise remain unquantified: where freshwater FIB data are actually collected and openly shared, and whether the cadence of sampling is fast enough to protect health.

These weaknesses matter because fecal contamination is a transient, weather-driven condition. Concentrations at a single site routinely vary by orders of magnitude over hours to days in response to rainfall, runoff, resuspension, and sunlight-driven inactivation \cite{whitman2004escherichia,whitman2004solar,traister2006variability}. A monitoring design that samples weekly, monthly, or quarterly, and often only during a summer bathing season, is therefore well suited to compliance reporting but poorly matched to the dynamics that determine acute risk on any given day.

Here we compile, harmonize, and analyze a global freshwater FIB database of 5{,}314{,}791 observations drawn from ten public sources and standardized to a common schema, indicator vocabulary, and unit. Using this dataset we (i) map the global distribution of freshwater FIB monitoring and show that openly shared data are concentrated in the United States, Europe, and Canada; (ii) characterize the indicator and ecosystem composition of the record; (iii) quantify per-site sampling cadence and show that only a small minority of sites are sampled at a frequency capable of resolving transient contamination; and (iv) use the rare high-frequency sites, including three contrasting case studies, to demonstrate that low-frequency sampling misrepresents safety in both transient and chronically polluted waters. We interpret the resulting geography as evidence that apparent data sparsity is, in substantial part, an open-data sharing barrier rather than an absence of monitoring, and we argue for both wider adoption of FAIR data practices \cite{wilkinson2016fair} and a shift toward higher-frequency, event-based, and predictive monitoring.

\section{Related Work}

\subsection{Global and harmonized water-quality data compilations}
A growing body of work aggregates water-quality observations across providers to enable large-sample analysis. In the United States, the Water Quality Portal unifies federal, state, and tribal records \cite{read2017water}, although reusing such multi-source data requires substantial harmonization and quality control \cite{sprague2017challenges,bousquin2024harmonizewq}. At continental to global scale, compilations such as AquaSat, GLORICH, the Global River Water Quality Archive, large-sample chemistry extensions of CAMELS, and global salinity syntheses have assembled physical and chemical water-quality variables for thousands of sites \cite{ross2019aquasat,hartmann2014brief,virro2021grqa,sterle2024camelschem,thorslund2020global}, and recent efforts have begun to argue for global microbial water-quality data infrastructure \cite{rose2023global,bain2014global}. These resources are invaluable, but they are typically organized around chemistry or single-nation coverage and do not characterize the spatial sharing structure or the temporal sampling cadence of fecal indicators specifically. We complement them with a freshwater-, FIB-focused harmonization across ten sources and an explicit analysis of those two dimensions.

\subsection{Fecal indicator bacteria, thresholds, and health risk}
The regulatory use of FIB rests on epidemiological studies relating indicator concentrations to swimming-associated illness \cite{cabelli1982swimmingassociated,prss1998review}, which underpin criteria such as single-sample and geometric-mean thresholds for \textit{E. coli} in fresh water \cite{wade2003environmental,wade2005rapidly,wade2010rapidly,boehm2008sea}. The choice of \textit{E. coli} as the freshwater indicator reflects its relationship to fecal contamination and health risk \cite{edberg2000escherichia,korajkic2018relationships}. A recurring theme in this literature is that thresholds are defined for compliance statistics and implicitly assume representative sampling. Our analysis quantifies how strongly that assumption is violated in practice and how the apparent safety of a water body depends on how often it is sampled.

\subsection{Temporal variability and high-frequency monitoring of FIB}
Numerous site studies document that fecal indicator concentrations vary sharply over short timescales, exhibiting diurnal cycles, rainfall- and runoff-driven spikes, and rapid decay \cite{traister2006variability,meays2006diurnal,desai2013escherichia,nguyen2022fecal}. This variability makes infrequent grab samples unreliable for classifying recreational safety \cite{whitman2004escherichia,kim2004public}, and it has motivated high-frequency sampling campaigns and statistical or data-driven nowcasting models that predict water quality from rainfall, turbidity, and discharge \cite{searcy2023highfrequency,searcy2021day,heasley2021systematic}. These approaches are demonstrated at a handful of intensively studied sites; what has been missing is a global accounting of how rare such high-frequency monitoring actually is and what it reveals when it exists. We provide that accounting and show that low-frequency sampling misleads in both transient and chronic regimes.

\subsection{Open environmental data and data-sharing barriers}
Finally, our interpretation of the global map connects to work on open environmental data. The FAIR principles articulate norms for findable, accessible, interoperable, and reusable data \cite{wilkinson2016fair}, yet a range of technical, institutional, and social barriers impede the open sharing of water-quality and other environmental monitoring data \cite{brasier2017barriers,sugg2021social,hazell2023sociotechnical,heacock2022enhancing}, and monitoring capacity itself varies systematically with national development \cite{kirschke2020capacity,kumpel2020data}. These dynamics imply that an apparent absence of data in global repositories often reflects where data are shared rather than where they are collected, a distinction our regional analysis makes explicit.

%Text here ===>>>


%%

%  Numbered lines in equations:
%  To add line numbers to lines in equations,
%  \begin{linenomath*}
%  \begin{equation}
%  \end{equation}
%  \end{linenomath*}



%% Enter Figures and Tables near as possible to where they are first mentioned:
%
% DO NOT USE \psfrag or \subfigure commands.
%
% Figure captions go below the figure.
% Acronyms used in figure captions will be spelled out in the final, published version.

% Table titles go above tables;  other caption information
%  should be placed in last line of the table, using
% \multicolumn2l{$^a$ This is a table note.}
% NOTE that there is no difference between table caption and table heading in the final, published version
%
%----------------
% EXAMPLE FIGURES
%
% \begin{figure}
% \includegraphics{example.png}
% \caption{caption}
% \end{figure}
%
% Giving latex a width will help it to scale the figure properly. A simple trick is to use \textwidth. Try this if large figures run off the side of the page.
% \begin{figure}
% \noindent\includegraphics[width=\textwidth]{anothersample.png}
%\caption{caption}
%\label{pngfiguresample}
%\end{figure}
%
%
% If you get an error about an unknown bounding box, try specifying the width and height of the figure with the natwidth and natheight options. This is common when trying to add a PDF figure without pdflatex.
% \begin{figure}
% \noindent\includegraphics[natwidth=800px,natheight=600px]{samplefigure.pdf}
%\caption{caption}
%\label{pdffiguresample}
%\end{figure}
%
%
% PDFLatex does not seem to be able to process EPS figures. You may want to try the epstopdf package.
%

%
% ---------------
% EXAMPLE TABLE
%
% \begin{table}
% \caption{Time of the Transition Between Phase 1 and Phase 2$^{a}$}
% \centering
% \begin{tabular}{l c}
% \hline
%  Run  & Time (min)  \\
% \hline
%   $l1$  & 260   \\
%   $l2$  & 300   \\
%   $l3$  & 340   \\
%   $h1$  & 270   \\
%   $h2$  & 250   \\
%   $h3$  & 380   \\
%   $r1$  & 370   \\
%   $r2$  & 390   \\
% \hline
% \multicolumn{2}{l}{$^{a}$Footnote text here.}
% \end{tabular}
% \end{table}

%%%%%%%%%%%%%%%%%%%%%%%%%%%%%%%%%%%%%%%%%%%%%%%
% SIDEWAYS FIGURES and TABLES
% AGU prefers the use of {sidewaystable} over {landscapetable} as it causes fewer problems.
%
% \begin{sidewaysfigure}
% \includegraphics[width=20pc]{figsamp}
% \caption{caption here}
% \label{newfig}
% \end{sidewaysfigure}
%
%  \begin{sidewaystable}
%  \caption{Caption here}
% \label{tab:signif_gap_clos}
%  \begin{tabular}{ccc}
% one&two&three\\
% four&five&six
%  \end{tabular}
%  \end{sidewaystable}

%% If using numbered lines, please surround equations with \begin{linenomath*}...\end{linenomath*}
%\begin{linenomath*}
%\begin{equation}
%y|{f} \sim g(m, \sigma),
%\end{equation}
%\end{linenomath*}

%%% End of body of article

%%%%%%%%%%%%%%%%%%%%%%%%%%%%%%%%%%%%%%%%%%%%%%%
%% Optional Appendices go here
%
% The \appendix command resets counters and redefines section heads
%
% After typing \appendix
%
%\section{Here Is Appendix Title}
% will show
% A: Here Is Appendix Title
%
%\appendix
%\section{Here is a sample appendix}

%%%%%%%%%%%%%%%%%%%%%%%%%%%%%%%%%%%%%%%%%%%%%%%
% Optional Glossary, Notation or Acronym section goes here:
%
% Glossary is only allowed in Reviews of Geophysics
%  \begin{glossary}
%  \term{Term}
%   Term Definition here
%  \term{Term}
%   Term Definition here
%  \term{Term}
%   Term Definition here
%  \end{glossary}


%%%%%%%%%%%%%%%%%%%%%%%%%%%%%%%%%%%%%%%%%%%%%%%
% Acronyms
%% NOTE that acronyms in the final published version will be spelled out when used in figure captions.
%   \begin{acronyms}
%   \acro{Acronym}
%   Definition here
%   \acro{EMOS}
%   Ensemble model output statistics
%   \acro{ECMWF}
%   Centre for Medium-Range Weather Forecasts
%   \end{acronyms}


%%%%%%%%%%%%%%%%%%%%%%%%%%%%%%%%%%%%%%%%%%%%%%%
% Notation
%   \begin{notation}
%   \notation{$a+b$} Notation Definition here
%   \notation{$e=mc^2$}
%   Equation in German-born physicist Albert Einstein's theory of special
%  relativity that showed that the increased relativistic mass ($m$) of a
%  body comes from the energy of motion of the body—that is, its kinetic
%  energy ($E$)—divided by the speed of light squared ($c^2$).
%   \end{notation}




%%%%%%%%%%%%%%%%%%%%%%%%%%%%%%%%%%%%%%%%%%%%%%%
%
% DATA SECTION and ACKNOWLEDGMENTS
%
%%%%%%%%%%%%%%%%%%%%%%%%%%%%%%%%%%%%%%%%%%%%%%%

\section*{Open Research Section}
This section MUST contain a statement that describes where the data supporting the conclusions can be obtained. Data cannot be listed as ''Available from authors'' or stored solely in supporting information. Citations to archived data should be included in your reference list. Wiley will publish it as a separate section on the paper’s page. Examples and complete information are here:
https://www.agu.org/Publish with AGU/Publish/Author Resources/Data for Authors

\section*{As Applicable – Inclusion in Global Research Statement}
The Authorship: Inclusion in Global Research policy aims to promote greater equity and transparency in research collaborations. AGU Publications encourage research collaborations between regions, countries, and communities and expect authors to include their local collaborators as co-authors when they meet the AGU Publications authorship criteria (described here: https://www.agu.org/publications/authors/policies\#authorship). Those who do not meet the criteria should be included in the Acknowledgement section. We encourage researchers to consider recommendations from The TRUST CODE - A Global Code of Conduct for Equitable Research Partnerships (https://www.globalcodeofconduct.org/) when conducting and reporting their research, as applicable, and encourage authors to include a disclosure statement pertaining to the ethical and scientific considerations of their research collaborations in an “Inclusion in Global Research Statement’ as a standalone section in the manuscript following the Conclusions section. This can include disclosure of permits, authorizations, permissions and/or any formal agreements with local communities or other authorities, additional acknowledgements of local help received, and/or description of end-users of the research. You can learn more about the policy in this editorial. Example statements can be found in the following published papers: 
Holt et al. (https://agupubs.onlinelibrary.wiley.com/doi/full/10.1029/2022JG007188), 
Sánchez-Gutiérrez et al. (https://agupubs.onlinelibrary.wiley.com/doi/abs/10.1029/2023JG007554), 
Tully et al. (https://agupubs.onlinelibrary.wiley.com/doi/epdf/10.1029/2022JG007128) 
Please note that these statements are titled as “Global Research Collaboration Statements” from a previous pilot requirement in JGR Biogeosciences. The pilot has ended and statements should now be titled “Inclusion in Global Research Statement”.

\section*{Conflict of Interest disclosure}
Please include a comprehensive conflict of interest statement that reflect all conflicts of interest for all involved authors. If there are no conflicts of interest please state, “The authors declare there are no conflicts of interest for this manuscript.”


\acknowledgments
Enter acknowledgments here. This section is to acknowledge funding, thank colleagues, enter any secondary affiliations, and so on.




%%%%%%%%%%%%%%%%%%%%%%%%%%%%%%%%%%%%%%%%%%%%%%%
% REFERENCES and BIBLIOGRAPHY
%
% \bibliography{<name of your .bib file>} don't specify the file extension
% don't specify bibliographystyle
%
%%%%%%%%%%%%%%%%%%%%%%%%%%%%%%%%%%%%%%%%%%%%%%%

\bibliography{refs}



%Reference citation instructions and examples:
%
% Please use ONLY \cite and \citeA for reference citations.
% \cite for parenthetical references
% ...as shown in recent studies (Simpson et al., 2019)
% \citeA for in-text citations
% ...Simpson et al. (2019) have shown...
%
%
%...as shown by \citeA{jskilby}.
%...as shown by \citeA{lewin76}, \citeA{carson86}, \citeA{bartoldy02}, and \citeA{rinaldi03}.
%...has been shown \cite{jskilbye}.
%...has been shown \cite{lewin76,carson86,bartoldy02,rinaldi03}.
%... \cite <i.e.>[]{lewin76,carson86,bartoldy02,rinaldi03}.
%...has been shown by \cite <e.g.,>[and others]{lewin76}.
%
% apacite uses < > for prenotes and [ ] for postnotes
% DO NOT use other cite commands (e.g., \citet, \citep, \citeyear, \nocite, \citealp, etc.).
%



\end{document}



More Information and Advice:

%%%%%%%%%%%%%%%%%%%%%%%%%%%%%%%%%%%%%%%%%%%%%%%
%
%  SECTION HEADS
%
%%%%%%%%%%%%%%%%%%%%%%%%%%%%%%%%%%%%%%%%%%%%%%%

% Capitalize the first letter of each word (except for
% prepositions, conjunctions, and articles that are
% three or fewer letters).

% AGU follows standard outline style; therefore, there cannot be a section 1 without
% a section 2, or a section 2.3.1 without a section 2.3.2.
% Please make sure your section numbers are balanced.
% ---------------
% Level 1 head
%
% Use the \section{} command to identify level 1 heads;
% type the appropriate head wording between the curly
% brackets, as shown below.
%
%An example:
%\section{Level 1 Head: Introduction}
%
% ---------------
% Level 2 head
%
% Use the \subsection{} command to identify level 2 heads.
%An example:
%\subsection{Level 2 Head}
%
% ---------------
% Level 3 head
%
% Use the \subsubsection{} command to identify level 3 heads
%An example:
%\subsubsection{Level 3 Head}
%
%---------------
% Level 4 head
%
% Use the \subsubsubsection{} command to identify level 3 heads
% An example:
%\subsubsubsection{Level 4 Head} An example.
%
%%%%%%%%%%%%%%%%%%%%%%%%%%%%%%%%%%%%%%%%%%%%%%%
%
%  IN-TEXT LISTS
%
%%%%%%%%%%%%%%%%%%%%%%%%%%%%%%%%%%%%%%%%%%%%%%%
%
% Do not use bulleted lists; enumerated lists are okay.
% \begin{enumerate}
% \item
% \item
% \item
% \end{enumerate}
%
%%%%%%%%%%%%%%%%%%%%%%%%%%%%%%%%%%%%%%%%%%%%%%%
%
%  EQUATIONS
%
%%%%%%%%%%%%%%%%%%%%%%%%%%%%%%%%%%%%%%%%%%%%%%%

% Single-line equations are centered.
% Equation arrays will appear left-aligned.

Math coded inside display math mode \[ ...\]
 will not be numbered, e.g.,:
 \[ x^2=y^2 + z^2\]

 Math coded inside \begin{equation} and \end{equation} will
 be automatically numbered, e.g.,:
 \begin{equation}
 x^2=y^2 + z^2
 \end{equation}


% To create multiline equations, use the
% \begin{eqnarray} and \end{eqnarray} environment
% as demonstrated below.
\begin{eqnarray}
  x_{1} & = & (x - x_{0}) \cos \Theta \nonumber \\
        && + (y - y_{0}) \sin \Theta  \nonumber \\
  y_{1} & = & -(x - x_{0}) \sin \Theta \nonumber \\
        && + (y - y_{0}) \cos \Theta.
\end{eqnarray}

%If you don't want an equation number, use the star form:
%\begin{eqnarray*}...\end{eqnarray*}

% Break each line at a sign of operation
% (+, -, etc.) if possible, with the sign of operation
% on the new line.

% Indent second and subsequent lines to align with
% the first character following the equal sign on the
% first line.

% Use an \hspace{} command to insert horizontal space
% into your equation if necessary. Place an appropriate
% unit of measure between the curly braces, e.g.
% \hspace{1in}; you may have to experiment to achieve
% the correct amount of space.


%%%%%%%%%%%%%%%%%%%%%%%%%%%%%%%%%%%%%%%%%%%%%%%
%
%  EQUATION NUMBERING: COUNTER
%
%%%%%%%%%%%%%%%%%%%%%%%%%%%%%%%%%%%%%%%%%%%%%%%

% You may change equation numbering by resetting
% the equation counter or by explicitly numbering
% an equation.

% To explicitly number an equation, type \eqnum{}
% (with the desired number between the brackets)
% after the \begin{equation} or \begin{eqnarray}
% command.  The \eqnum{} command will affect only
% the equation it appears with; LaTeX will number
% any equations appearing later in the manuscript
% according to the equation counter.
%

% If you have a multiline equation that needs only
% one equation number, use a \nonumber command in
% front of the double backslashes (\\) as shown in
% the multiline equation above.

% If you are using line numbers, remember to surround
% equations with \begin{linenomath*}...\end{linenomath*}

%  To add line numbers to lines in equations:
%  \begin{linenomath*}
%  \begin{equation}
%  \end{equation}
%  \end{linenomath*}



